% ==================== 5. Conclusion ====================
\section{Conclusion}
\label{sec:concl}

We revisited what a real-pair SIRR protocol silently assumes, for the practical case of
automotive tinted glass. The occlusion region is real, it is measurable, and its dominant
effect is dominated by the reflection deficit ($16.8\%$) rather than by target darkening ($3.7\%$),
while the input still falls $13\%$ below its own target inside $\Omega$ --- which is why a globally
strength-calibrated correction over-corrects inside it. More usefully for practitioners, our
controlled comparison shows that the benefit of an occlusion-mask loss does \emph{not} come
from the mask: an area-matched random region is at least as good, and the active ingredient is
the normalization of the error by the local reflection amplitude. Where the signal is
demonstrably present --- the reflection layer itself --- replacing a gradient-comparison
pseudo-label by true supervision triples the usable mask quality for free. We also give a
closed-form bound that tells, before running anything, whether a region- or boundary-scale
term can be validated by a global score at all, and an interventional test that exposes
photometric conditions which do not actually condition. All negative results, the dataset and
the evaluation protocol are released so that the next attempt can start from the right
question.

\section*{Reproducibility note}
This submission extends an anonymized project report; per-arm logs, the frozen split hash and
the evaluation script are included as supplementary material. Three-seed arms use seeds
101/202/303.
